\begin{figure*}[t]
  \centering
  \includegraphics[width=\textwidth]{pipeline.pdf}
  \caption{FILA-Net pipeline. Top: inference on one image. Bottom: training,
  validation and calibration, all scored with the official PQ.}
  \label{fig:pipeline}
\end{figure*}

\paragraph{Data and splits.}
The training set has 707 observations, 1{,}154 annotator label sets and
8{,}199 filaments (411 observations labeled once, 145 twice, 151 three
times). Every label set is a training sample. Observations of the same UT
day show the same Sun from different GONG sites minutes apart (101 training
observations share their day), so we use 5 folds grouped by day: all
observations of a day and all their label sets belong to one fold, and the
OOF score is not inflated by near-duplicates. The assignment is frozen in
the repository (\texttt{configs/folds\_day5.csv}). Only the annotations
provided by the competition are used.

\paragraph{Inputs.}
Each grayscale image becomes four channels: intensity normalized between
its 1st and 99th percentiles, CLAHE~\cite{zuiderveld1994clahe}, a ridge
response (difference of Gaussians, $\sigma=1.2$ and $5$) that highlights
thin dark threads, and the normalized radial position, which lets the
network account for limb darkening (Fig.~\ref{fig:inputs}).

\begin{figure*}[t]
  \centering
  \includegraphics[width=\textwidth]{inputs_and_annotations.png}
  \caption{The four input channels of a training observation and the
  filaments drawn by its three annotators (18, 18 and 22 filaments): the
  label sets disagree on which faint features are filaments and where they
  end, which bounds the achievable PQ.}
  \label{fig:inputs}
\end{figure*}

\paragraph{Network.}
An ImageNet-initialized ConvNeXt encoder~\cite{liu2022convnext} (via
timm~\cite{rw2019timm}) feeds a feature pyramid~\cite{lin2017fpn} fused at
stride~4 with 128 channels, refined by a multi-dilation context block, a
high-frequency branch and a limb-aware gate, and upsampled to full
resolution. $1\times1$ heads predict the region, the centerline (dilated
skeleton), the boundary, the per-filament normalized distance, the axial
orientation $(\cos2\theta,\sin2\theta)$ of the skeleton, and auxiliary
width, curvature, endpoint, junction and aleatoric-uncertainty
maps~\cite{kendall2017uncertainties}. All targets are derived from the
segmentation polygons only. We train two encoders, ConvNeXt-T
(\texttt{fb\_in1k}) and ConvNeXt-S (\texttt{fb\_in22k\_ft\_in1k}), on the
same folds.

\paragraph{Training.}
The loss combines BCE and Dice on region (0.25, 0.30) and centerline (0.10,
0.15), BCE on the boundary (0.08), a smooth-$L_1$ distance term (0.12),
clDice~\cite{shit2021cldice} for topological continuity (0.05), a cosine
loss on the orientation (0.08) and small weights ($\le0.05$) on the
auxiliary heads. We train on random $1024\times1024$ crops (70\% centered on
a filament) with flips, AdamW (lr $2\cdot10^{-4}$, one warm-up epoch, cosine
decay to $10^{-6}$, weight decay $10^{-4}$), batch 2, mixed precision and
gradient clipping, for 30 epochs. Every two epochs the held-out fold is
predicted in full and scored with the official PQ, and the best epoch is
kept (epochs 22--30; Fig.~\ref{fig:curves}).

\paragraph{Inference.}
Images are processed in $1024$-pixel tiles with $256$-pixel overlap blended
by a Hann window, with test-time augmentation by horizontal and vertical
flips (the orientation field is transformed consistently), and the maps of
the ten fold models (2 encoders $\times$ 5 folds) are averaged.

\paragraph{Instance reconstruction.}
The region probability is thresholded and small components are removed; the
thresholded centerline is skeletonized to obtain seeds. Skeleton endpoints of
different fragments are joined when the gap is short, aligned with the
predicted orientation on both sides, supported by the region probability,
free of boundary evidence and compatible in width (greedy acceptance in order
of cost, without cycles). A marker-controlled
watershed~\cite{beucher1993watershed} on the distance map, penalized at
predicted boundaries, then assigns every foreground pixel to one filament,
giving disjoint instances. Its 14 parameters (thresholds, minimum areas,
repair distances and angles, a minimum mean region probability per instance)
are calibrated on OOF predictions (Section~\ref{sec:evaluation}). Masks are
written as COCO RLE~\cite{lin2014coco}.

\paragraph{Learned filament filter.}
Predictions are disjoint, so a ground-truth filament is hit by at most one
of them and removing one prediction changes nothing else. Let $a_j$ be the
sum of the IoUs of the matches of prediction $j$ over the label sets of its
image and $d_j$ the sum, over these label sets, of $1/2$ if $j$ matches
nothing and $h/2$ if it matches $h$ filaments. For a kept set $K$ and $G$
ground-truth filaments,
\[
  \mathrm{PQ}(K)=\frac{\sum_{j\in K}a_j}{G/2+\sum_{j\in K}d_j},
\]
exactly (checked against the official counts on all OOF images), so $j$
should be kept iff $a_j/d_j>\mathrm{PQ}$. We regress
$y_j=a_j/(2d_j)$ with a histogram gradient-boosted regressor
\cite{ke2017lightgbm,pedregosa2011sklearn} (300 trees, 15 leaves) on 32
features of each instance: 25 from the predicted maps (area, elongation,
skeleton length, mean/max probabilities, contrast with a 3-pixel ring,
radial position, rank and number of instances in the image) and 7 from the
H-$\alpha$ image itself (intensity, contrast with its surroundings, ridge
response). Instances with predicted $y$ below a threshold are dropped. Two
candidate sets are compared: the calibrated instances, and a permissive set
(no minimum probability, area $\ge100$ pixels) that the filter thins out.
The filter is cross-fitted on OOF predictions and uses only the image being
segmented.
